% =====================================================================
% CHAPTER 7: Societal, Environmental Impact and Sustainability
% Restructured 2026-07-31 per external examiner's final feedback:
%   7.1 Introduction
%   7.2 Impact of the Thesis on Societal, Health, and Cultural Issues
%   7.3 Safety and Ethical Considerations
%   7.4 Alignment with Sustainable Development Goals (SDGs)
%   7.5 Summary
% Environmental and sustainability material folded into the SDG section
% (SDG 12 in particular). Chapter label ch:socio kept (referenced by
% ch1 organization list and app_i PO tracker). POs demonstrated:
% PO6 (engineer and society), PO7 (environment and sustainability).
% =====================================================================
\chapter{Societal, Environmental Impact and Sustainability}
\label{ch:socio}

\section{Introduction}
The value of this work is not exhausted by its benchmark numbers. A depth
sensor that runs its whole computation on the device, emits nothing into
the scene, and is built from a handful of commodity parts carries
consequences for the people around it, for the environment it operates
in, and for the wider goals that engineering is meant to serve. This
chapter draws those consequences out. Section~\ref{sec:soc_impact}
examines the societal, health, and cultural effects that follow from
where the computation happens and how the scene is sensed.
Section~\ref{sec:safety_ethics} turns to safety and to the ethical
questions a camera-carrying, autonomy-enabling system raises.
Section~\ref{sec:sdg} maps the work onto three of the United Nations
Sustainable Development Goals and, in doing so, sets out its
environmental and sustainability case. Section~\ref{sec:soc_summary}
draws the chapter together.

\FloatBarrier
\section{Impact of the Thesis on Societal, Health, and Cultural Issues}
\label{sec:soc_impact}
The most direct societal consequence of this work follows from where the
computation happens. Because the network runs the depth computation on
the device rather than in the cloud, the video stream stays local: no
imagery of the people or spaces that a robot, drone, or headset observes
ever needs to leave the sensor. This matters for privacy in a way that a
cloud-processed depth pipeline cannot match, since a system that
transmits continuous video to a remote server creates, by construction, a
stream of personal imagery that must then be secured, retained, and
governed. Keeping the raw imagery on the device removes that stream at
the source. The same property eases the legal and cultural concerns that
accompany camera-carrying machines. Regulations on personal-data
transfer, and the reasonable discomfort of people who share a space with
an observing device, both attach primarily to imagery that is recorded
or transmitted; an embedded sensor that consumes its own video internally
and emits only geometry is easier to deploy responsibly in homes,
workplaces, and public settings.

The choice of sensing modality carries a health implication that is easy
to overlook. A passive stereo rig emits no laser and no structured
infrared light; it only observes. Bystanders are therefore not exposed to
any emitted radiation, a real distinction from active depth sensors,
whose emitters must be engineered and certified for eye safety,
particularly around children and in crowded indoor spaces. A depth sensor
that adds no new emission to a shared room is, by that measure, gentler
on the people in it.

Finally, the work has an accessibility dimension with a cultural reach.
A compact, low-cost, passive depth sensor lowers the barrier to robotics
and assistive technology in settings where specialized depth hardware is
unaffordable. Depth perception built from two commodity image sensors and
an embedded accelerator puts spatial sensing within reach of student
projects, small workshops, assistive devices, and products aimed at
price-sensitive markets, communities that a sensor costing several times
as much would simply exclude. Widening who can build and afford spatially
aware machines is itself a societal outcome, and it follows directly from
the design constraints this thesis set out to satisfy.

\FloatBarrier
\section{Safety and Ethical Considerations}
\label{sec:safety_ethics}
Safety is the first concern for any system that senses depth so that a
machine can move. A mobile robot or drone acts safely in
three-dimensional space only to the extent that its depth perception is
reliable, and the evaluation of Chapter~\ref{ch:results} was built around
exactly the failure modes that make depth unreliable in the field: the
model was tested not only in its training domain but zero-shot on unseen
real scenes, under controlled rectification error, and on the physical
stereo rig itself. On-device inference also removes the network round
trip from the control loop, so the depth estimate a robot acts on does
not depend on the availability or latency of a wireless link; a platform
whose obstacle avoidance survives a lost connection is categorically safer
than one whose perception stalls with the network. The passive rig
reinforces this: with no laser and no projected pattern, the sensor poses
no eye-safety hazard to the people it works alongside.

Reliability, though, has limits that honesty requires stating. The
network is trained on synthetic data and is not perfect on real scenes;
its errors concentrate at depth boundaries and on thin, repeated, or
reflective structures, and its accuracy degrades gracefully rather than
failing silently, but it degrades. A system that acts on this depth
should therefore treat it as one cue among several rather than as ground
truth, and this thesis reports the error rates plainly, including where
they are worst, so that a downstream integrator can design for them
instead of being surprised by them.

The ethical questions extend past safety. A device that perceives
three-dimensional space is inherently dual-use: the same depth stream
that lets a robot avoid a person can, in another deployment, help track
one. The design choices here lean against the harmful use, on-device
processing keeps imagery from being collected centrally, and passive
sensing avoids the covert illumination that active depth can use, but no
architecture can prevent misuse on its own, and it would be dishonest to
claim otherwise. Two further points of engineering ethics bear on this
work directly. First, the data: the real pairs used for qualification
were captured by the authors in their own institutional spaces, with no
third parties as subjects, and no personal imagery was collected,
retained, or distributed. Second, the reporting: every accuracy, latency,
and efficiency figure in this thesis is tied to a specific scripted run
whose configuration is logged, comparisons against published methods use
each paper's own reported numbers, and the model's shortcomings are
stated alongside its strengths, so that the record is one a reader can
trust and reproduce rather than one tuned to flatter.

\FloatBarrier
\section{Alignment with Sustainable Development Goals (SDGs)}
\label{sec:sdg}
The work aligns with three of the United Nations Sustainable Development
Goals. Each connection is concrete rather than aspirational, and the
environmental case for the design is made here, under the goal it serves
most directly.

\paragraph{SDG 9: Industry, Innovation and Infrastructure.}
This goal calls for resilient infrastructure and for broad, affordable
access to the fruits of innovation. A stereo network that delivers usable
depth within the memory, latency, and power budget of a low-cost edge
board moves a capability that until recently required either expensive
sensors or a datacenter onto hardware that small manufacturers,
researchers, and startups can actually afford. By showing that a design
methodology, controlled ablation under a hard budget, can produce a
competitive model at a few million parameters, the thesis also
contributes a repeatable way of building efficient perception rather than
a single artifact, which is the kind of transferable innovation this goal
seeks.

\paragraph{SDG 11: Sustainable Cities and Communities.}
This goal includes safe, accessible, and sustainable mobility and living
spaces. Reliable, low-cost, on-device depth perception is an enabling
technology for exactly the machines that operate in shared human
environments: service and delivery robots, inspection drones, and
assistive devices for people with limited mobility or vision. Because the
sensor keeps imagery local and emits nothing, it can be deployed in homes,
sidewalks, and public buildings with a lighter privacy and safety
footprint than a cloud-connected or actively illuminating alternative,
which is what makes such deployments socially sustainable rather than
merely technically possible.

\paragraph{SDG 12: Responsible Consumption and Production.}
This goal is where the environmental case sits, and it applies to the
hardware, the running energy, and the development of the model alike. On
the hardware side, a stereo rig uses two commodity image sensors rather
than a specialized time-of-flight or LiDAR unit and carries no emitter
hardware at all: no laser diode, no pattern projector, and none of the
associated drive electronics. Those commodity sensors are manufactured at
enormous scale for the phone and webcam markets, so a depth sensor
assembled from them draws on an established supply chain rather than
specialized parts, and its fewer, more common components make it easier
to source, repair, and replace over its service life. On the energy side,
the deployed network runs within a few watts to a few tens of watts and
computes each frame locally, avoiding the recurring per-frame cost of
transmitting video to a datacenter and back; because that transport cost
is paid for every frame over a device's whole life, removing it compounds
over the deployed fleet. On the development side, the production model was
trained in roughly one day on a single cloud GPU, and each ablation and
evaluation run cost a fraction of a dollar of cloud time, a footprint
kept deliberately small by the cheap controlled-ablation harness of
Chapter~\ref{ch:design}, which filtered design candidates before any
expensive full training run. Compute, in other words, was spent once, on
a design already validated by a cheaper process, which is responsible
production applied to model building itself.

\FloatBarrier
\section{Summary}
\label{sec:soc_summary}
This chapter traced the design properties of the system to their human
and environmental consequences. On-device inference keeps imagery local,
which protects the privacy of observed people and eases legal and
cultural concerns; passive sensing emits no laser and poses no eye-safety
hazard; and the low cost of the sensor widens access to assistive and
robotic technology. On safety and ethics, connection-independent depth
makes robots and drones safer, the thesis reports its errors honestly
rather than only its successes, and it acknowledges the dual-use nature
of depth sensing while leaning its design choices against misuse.
Finally, the work aligns with three Sustainable Development Goals,
affordable innovation (SDG 9), safe machines in shared spaces (SDG 11),
and responsible production (SDG 12), and it is under the last of these
that the environmental case is made: commodity parts over specialized
emitters, local computation over a repeated cloud round trip, and a
training footprint of about one GPU-day. These impacts follow directly
from the edge-deployment constraints that shaped the design.

% SOURCES: mainmatter2/ch7.tex prior version (societal/privacy,
%   legal/cultural, accessibility, safety, environment, and training-
%   footprint prose all carried over and re-sectioned). New prose: the
%   ethical-considerations paragraphs (dual-use, data provenance, honest
%   reporting) and the three SDG mappings.
% DROPPED: none substantive; the old separate "Environment and
%   Sustainability" section content was folded into SDG 12.
% MOVED-AWAY: none. Chapter label ch:socio retained.
